% Options for packages loaded elsewhere
\PassOptionsToPackage{unicode}{hyperref}
\PassOptionsToPackage{hyphens}{url}
\PassOptionsToPackage{dvipsnames,svgnames,x11names}{xcolor}
\documentclass[
]{article}
\usepackage{xcolor}
\usepackage{amsmath,amssymb}
\setcounter{secnumdepth}{5}
\usepackage{iftex}
\ifPDFTeX
  \usepackage[T1]{fontenc}
  \usepackage[utf8]{inputenc}
  \usepackage{textcomp} % provide euro and other symbols
\else % if luatex or xetex
  \usepackage{unicode-math} % this also loads fontspec
  \defaultfontfeatures{Scale=MatchLowercase}
  \defaultfontfeatures[\rmfamily]{Ligatures=TeX,Scale=1}
\fi
\usepackage{lmodern}
\ifPDFTeX\else
  % xetex/luatex font selection
\fi
% Use upquote if available, for straight quotes in verbatim environments
\IfFileExists{upquote.sty}{\usepackage{upquote}}{}
\IfFileExists{microtype.sty}{% use microtype if available
  \usepackage[]{microtype}
  \UseMicrotypeSet[protrusion]{basicmath} % disable protrusion for tt fonts
}{}
\makeatletter
\@ifundefined{KOMAClassName}{% if non-KOMA class
  \IfFileExists{parskip.sty}{%
    \usepackage{parskip}
  }{% else
    \setlength{\parindent}{0pt}
    \setlength{\parskip}{6pt plus 2pt minus 1pt}}
}{% if KOMA class
  \KOMAoptions{parskip=half}}
\makeatother
\usepackage{longtable,booktabs,array}
\usepackage{caption}
\captionsetup[table]{skip=6pt}
\newcounter{none} % for unnumbered tables
\usepackage{calc} % for calculating minipage widths
% Correct order of tables after \paragraph or \subparagraph
\usepackage{etoolbox}
\makeatletter
\patchcmd\longtable{\par}{\if@noskipsec\mbox{}\fi\par}{}{}
\makeatother
% Allow footnotes in longtable head/foot
\IfFileExists{footnotehyper.sty}{\usepackage{footnotehyper}}{\usepackage{footnote}}
\makesavenoteenv{longtable}
\setlength{\emergencystretch}{3em} % prevent overfull lines
\providecommand{\tightlist}{%
  \setlength{\itemsep}{0pt}\setlength{\parskip}{0pt}}
%% ============================================================
%% paper/preamble.tex
%% Injected into pandoc's default LaTeX template via
%% `pandoc --include-in-header`. The committed report.tex inlines
%% this content, so the PDF rebuilds with tectonic alone (no pandoc).
%% ============================================================

% arXiv single-column preprint style (George Kour, based on the NeurIPS
% style). Vendored as paper/arxiv.sty so the build needs no network for it.
\usepackage{arxiv}

% arxiv.sty sets \rmdefault to ptm (8-bit Adobe Times), which has no Unicode
% (TU) encoding under the XeTeX engine tectonic uses. Every shape then falls
% back, and crucially `TU/ptm/b/n' resolves to `TU/ptm/m/n' -- bold silently
% becomes regular, so section and subsection headings render at body weight.
% TeX Gyre Termes is the Times clone with real OpenType weights; fontspec is
% already loaded by unicode-math in pandoc's template.
\setmainfont{texgyretermes}[
  Extension      = .otf,
  UprightFont    = *-regular,
  BoldFont       = *-bold,
  ItalicFont     = *-italic,
  BoldItalicFont = *-bolditalic,
]

% arxiv.sty stamps "A Preprint" under the title and in the running header.
% Drop both: this is a technical report, not a preprint of a submission.
% \undertitle is emptied rather than removed so \maketitle's spacing still works.
\renewcommand{\undertitle}{}
\renewcommand{\headeright}{}

% --- Bibliography engine -----------------------------------------------------
% natbib (numeric) is wired in. report-filters.lua turns the 8 related-work
% arXiv links into \citep{...} while keeping the visible work name, and attaches
% \citep{ross2011dagger} at the "DAgger-style relabel" anchor. The References
% section (references.tex) therefore lists exactly the 9 cited works.
%
% To cite more of the text later: map another inline link (or phrase) to a bib
% key the same way in report-filters.lua — no edits to REPORT.md required.
\usepackage[numbers,sort&compress]{natbib}

% --- Unicode glyphs missing from the Times (ptm) text font under the XeTeX
%     engine tectonic uses. Map them to LaTeX equivalents so none silently
%     drop out (arrows and the U+2212 minus in particular). --------------------
\usepackage{newunicodechar}
\newunicodechar{→}{\ensuremath{\rightarrow}}
\newunicodechar{—}{---}
\newunicodechar{–}{--}
\newunicodechar{…}{\ldots}
\newunicodechar{−}{\ensuremath{-}}

% --- Keep the prose-heavy tables (practitioner guide, related work, checker
%     grid) inside the text block. pandoc emits proportional wrapping p-columns
%     for wide tables; a smaller table font plus tighter padding guarantees the
%     page body never overflows. -----------------------------------------------
\usepackage{etoolbox}
\AtBeginEnvironment{longtable}{\small}
\setlength{\tabcolsep}{4pt}
\setlength{\LTleft}{0pt}     % flush-left longtables (no unwanted centering gap)
\setlength{\LTright}{0pt}

% Let long URLs in hyperlinks break rather than overflow the margin.
\PassOptionsToPackage{hyphens}{url}
\usepackage{bookmark}
\IfFileExists{xurl.sty}{\usepackage{xurl}}{} % add URL line breaks if available
\urlstyle{same}
% fallback for those not using the hyperref driver hyperxmp:
\makeatletter
\@ifundefined{xmpquote}{\newcommand{\xmpquote}[1]{#1}}{}
\makeatother
\hypersetup{
  pdftitle={When Do Language Servers Help Coding Agents?},
  pdfauthor={Ian Barber},
  colorlinks=true,
  linkcolor={blue},
  filecolor={Maroon},
  citecolor={blue},
  urlcolor={blue},
  pdfcreator={LaTeX via pandoc}}

\title{When Do Language Servers Help Coding Agents?}
\author{Ian Barber}
\date{July 2026}

\begin{document}
\maketitle
\begin{abstract}
Coding agents are increasingly given language-server tooling on the
assumption that better code intelligence makes a better agent. We
explored that assumption against a capable baseline of grep, ranged
reads and a shell, across 7 models. Making LSP tools available did not
change anything: agents largely declined to use the operations they were
offered. Gains required a prompt, training, or a gate that blocked the
work. Go-to-definition helped when resolving the receiver's type
required a retrieval step rather than being already in context, and only
when the agent was instructed to use it. Definition spans saved tokens
only when they replaced the file read, which is workload dependent: on
one suite models reread on 35 of 36 injected spans. Merely adding
prompting to say the span was complete removed reads in just 2 cases:
adjusting the behavior required fine-tuning. On 12 seeded defects, a
type checker gating submission took accepted correct outcomes from 1 to
11. Earlier delivery of type errors made no difference, and actively
guiding the model to respond to earlier diagnostics performed worse than
no feedback at all.

\textbf{If you use a coding agent:} run your type checker at the end of
the turn as a blocking gate; add navigation tools only as the codebase
gets more complex, with an instruction to use them, and expect the most
gains when resolving a type means reaching into a stub, a generated
client or a compiled boundary; expect to save tokens only once the agent
stops reading the file, so measure that before counting the savings.
\end{abstract}

\section{Introduction}\label{introduction}

A language server answers questions about code. An agent with grep,
ranged reads and a shell already answers most of them itself, so what a
server adds depends on whether it answers a question the agent cannot
(IF), whether its answer is cheaper in tokens (FORM), and whether it
arrives at a better moment (WHEN). There is significant prior work, but
it is rarely measured against an agent that already reads code well;
that is the baseline we used throughout.

\textbf{IF.} Typed Holes~\citep{blinn2024typedholes} and
LSPRAG~\citep{go2025lsprag} push types and definitions into context on
the premise the information is missing.

\textbf{FORM.} CodeStruct~\citep{kim2026codestruct} argues structured
reads and edits use fewer tokens than whole-file handling.

\textbf{WHEN.} STALL+~\citep{liu2024stallplus} finds static-analysis
value varying with integration phase; CoCoGen~\citep{bi2024cocogen}
checks a coherent draft before retrieving context to repair it.

\section{Method}\label{method}

Most experiments ran in a custom harness: a fixed loop of six actions
over a Python workspace (grep, ranged read, whole-file read, line edit,
run tests, submit), and each experimental arm added or withheld one
semantic operation. The definition operation in the retrieval
experiments was a static AST resolver validated for equivalence against
Pyrefly. Live Pyrefly served the dispatch lookup arms directly. A
confirmation study ran
mini-swe-agent~\citep{minisweagent,yang2024sweagent} with Claude Sonnet
4.5 on 3 SWE-bench~\citep{jimenez2024swebench} tasks (sympy, astropy,
sphinx).

Note: We did not use an MCP server, skill or editor plugin.

{\def\LTcaptype{none} % do not increment counter
\begin{longtable}[]{@{}
  >{\raggedright\arraybackslash}p{(\linewidth - 8\tabcolsep) * \real{0.2000}}
  >{\raggedright\arraybackslash}p{(\linewidth - 8\tabcolsep) * \real{0.2000}}
  >{\raggedright\arraybackslash}p{(\linewidth - 8\tabcolsep) * \real{0.2000}}
  >{\raggedright\arraybackslash}p{(\linewidth - 8\tabcolsep) * \real{0.2000}}
  >{\raggedright\arraybackslash}p{(\linewidth - 8\tabcolsep) * \real{0.2000}}@{}}
\toprule\noalign{}
\begin{minipage}[b]{\linewidth}\raggedright
Suite
\end{minipage} & \begin{minipage}[b]{\linewidth}\raggedright
N
\end{minipage} & \begin{minipage}[b]{\linewidth}\raggedright
Source
\end{minipage} & \begin{minipage}[b]{\linewidth}\raggedright
Temp x seeds
\end{minipage} & \begin{minipage}[b]{\linewidth}\raggedright
Isolates
\end{minipage} \\
\midrule\noalign{}
\endhead
\bottomrule\noalign{}
\endlastfoot
Dispatch ladder & 15 tasks x 3 variants & Synthetic; typed receiver,
\textasciitilde10 same-named overrides (one buggy); type moves from
call-site annotation to construction site to factory indirection & 0 x 1
& Lookup when the type is already in context \\
Hidden-source dispatch & 15 tasks, 450 rollouts & Same suite, with
application and test source left on disk but not pasted into the prompt
and the receiver construction redacted in the quoted asserts & 0.7 x 5 &
Lookup when the type costs a retrieval step \\
Whole-file baseline & 11 tasks x 4 seeds x 3 models & Constructed misuse
over vendored toolz 1.1.0 and more-itertools 11.1.0 & 0.7 x 4 & Span
cost against a whole-file read \\
Text-retrieval baseline & 11 tasks & As above & 0 x 1 & Span cost
against grep with ranged reads \\
Definition spans & 12 instances per arm & Synthetic arithmetic
overrides; the delivered span is mechanically verified to contain the
fix & 0 x 1 & Whether the span substitutes for the read \\
Delivery timing & 14 tasks x 6 arms, 168 rollouts per arm & Synthetic
single-file authoring; no feedback, 2 batched and 3 live delivery points
& 0.7 x 12 & When a diagnostic should arrive \\
Checker grid & 12 defect/clean pairs & Revision task: model reviews a
pre-written draft carrying a seeded defect that passes the visible test;
clean controls are validated gold & 0 x 1 & Timing of one diagnostic \\
Shell-agent case study & 3 SWE-bench tasks & sympy, astropy and sphinx,
through off-the-shelf mini-swe-agent with Claude Sonnet 4.5 and an
AST-backed definition command & one seed, 60-step cap & Whether the
constructed findings show up in a real agent on real repositories \\
\end{longtable}
}

Local models were Qwen2.5-Coder-7B, Qwen3.5-27B and Qwen3.6-27B; Claude
Sonnet 4.5, DeepSeek v3.1, GLM-4.6 and GPT-5.6 ``Luna'' ran through an
API. Most runs were temperature 0.7 with repeated seeds; some held-out
retests ran at temperature 0. Workspaces are small enough to read any
file completely; we did not test dynamic behavior, wrong annotations and
\texttt{Any}-heavy boundaries. Code and artifacts:
https://github.com/ianbarber/lsps-for-llms.

Measures: * \textbf{Resolution}: fixed the intended target and passed a
held-out test written against the specification, separate from the
visible one. * \textbf{Cost}: total tokens, input plus output, over the
whole trajectory. * \textbf{Election}: how often the agent chose an
offered semantic operation. * \textbf{Substitution}: reads of the
defining file after a span was delivered.

Checker experiments also report wall time. Measurements are per-task;
rollouts are grouped per task to get observations for intervals.

\section{Results}\label{results}

\subsection{IF: does delivering the semantic information change what the
agent
finds?}\label{if-does-delivering-the-semantic-information-change-what-the-agent-finds}

In the \emph{dispatch ladder} suite all source was provided in the
initial prompt, and the variants differed only in how far the type sits
from the call site: on it, at the construction site, or behind a
factory. The agent grepped 0.3 times per task and re-read source it had
already been given in the prompt, navigating to the buggy override on
the types available to it: text search and defn resolved equally at
every level.

For tokens, cost was flat too: 1.03x text search when lookup was merely
available, 0.96x when prompted, 0.94x to 1.06x across variants. The
difference was substitution: the prompted agent swapped the span for its
file read, resulting in 0.07 whole-file reads per task against 1.00,
while the unprompted agent made both calls (0.80 lookups on top of 1.00
reads).

In the \emph{hidden-source dispatch} suite the tasks are similar, but
the source is not provided in the prompt. To avoid an initial whole-file
read the use-site line and column are provided, but the type is no
longer free in context and has to be fetched by read, grep or lookup.
This suite ran three arms: text search alone (grep\_base), lookup
available but unprompted (defn\_avail), and lookup with an instruction
to use it (defn\_prompt).

{\def\LTcaptype{none} % do not increment counter
\begin{longtable}[]{@{}lrlrr@{}}
\toprule\noalign{}
Arm & Resolved & Difference vs grep & Greps & Whole-file reads \\
\midrule\noalign{}
\endhead
\bottomrule\noalign{}
\endlastfoot
text search & 70/75 & --- & 1.52 & 2.17 \\
lookup available & 71/75 & +0.013 {[}−0.040, +0.067{]} & 1.47 & 2.04 \\
lookup prompted & 75/75 & +0.067 {[}+0.013, +0.133{]} & 0.36 & 0.68 \\
\end{longtable}
}

\emph{Per-task resolution over 5 seeds at temperature 0.7; 95\% task
bootstrap intervals over 15 tasks; greps and reads are per-task means.
In the matched visible control every arm resolved all 75 rollouts.}

Availability raised resolution on 3 but lowered it on 2 tasks. Prompting
raised the resolution rate on 4 tasks with no impact on others. Tokens
barely moved (0.983 prompted, 1.005 available, on tasks both arms
solved).

We also verified that types are beneficial regardless of the tool as
part of the dispatch ladder. Removing annotations hurt accurate
retrieval, while the text agent succeeded 14-15/15 if the annotations
were present at call site, construction site, or behind indirection.

\subsection{FORM: does a compact span cost less than reading the
file?}\label{form-does-a-compact-span-cost-less-than-reading-the-file}

In the \emph{whole-file baseline} suite, the whole-file arm cost 3.35x
the definition arm in total tokens for Qwen3.6-27B, 3.49x for Sonnet 4.5
and 4.12x for DeepSeek v3.1. A capable agent rarely reads a whole file
though.

The \emph{text-retrieval baseline} suite explored the same 11 tasks of
real library source. Qwen3.5-27B spent 1,602 mean tokens with text
retrieval and 1,235 with definitions, a paired ratio of 1.297 {[}1.093,
1.527{]}. Definitions were cheaper on 10 of 11 and both arms solved each
task.

In order to benefit from this efficiency gain, we need to establish
whether the tool call augmented or substituted existing file reads. This
behavior appears to be workload-dependent: in the \emph{definition
spans} suite, injected defn spans were followed by a read of the
defining file on 35 of 36 instances. Prompting that the span was
complete removed just 2 of those reads, while adding a per-task cost of
294 tokens.

{\def\LTcaptype{none} % do not increment counter
\begin{longtable}[]{@{}lrr@{}}
\toprule\noalign{}
Model & Reread, injected span & Reread, plus sufficiency instruction \\
\midrule\noalign{}
\endhead
\bottomrule\noalign{}
\endlastfoot
Qwen3.6-27B & 11/12 & 12/12 \\
Sonnet 4.5 & 12/12 & 12/12 \\
DeepSeek v3.1 & 12/12 & 10/12 \\
\end{longtable}
}

The \emph{shell-agent case study} showed similar results on real
repositories: a manual cross-check matched 16 of 18 definition calls to
a later read of the file just resolved.

Training can impact substitution. A separate DAgger-style
relabel~\citep{ross2011dagger} on 39 demonstrations, none of which
required a read, cut the reread from 11 of 12 to 0 of 12 on held-out
instances with different seeds and templates. Mean input fell from 1,157
to 748 tokens, and from 1,493 to 938 on the nine tasks both arms solved,
a 1.59x saving. Held-out pass slipped from 11 to 10 of 12, one rescue
against two losses, which at 12 instances and one seed is
indistinguishable from noise. We held back a further 12 instances,
unused until the end: there the reread went from 12 of 12 to 1 of 12,
mean input from 1,223 to 729 tokens, a 1.72x saving on the eleven both
arms passed, and held-out pass from 12 to 11 of 12.

We also used the \emph{definition spans} suite to test sufficiency:
whether the model reads after a defn span is delivered if the
information contained is insufficient. A model that did not read would
gain token efficiency at the cost of correctness.

{\def\LTcaptype{none} % do not increment counter
\begin{longtable}[]{@{}
  >{\raggedright\arraybackslash}p{(\linewidth - 6\tabcolsep) * \real{0.2000}}
  >{\raggedleft\arraybackslash}p{(\linewidth - 6\tabcolsep) * \real{0.2667}}
  >{\raggedleft\arraybackslash}p{(\linewidth - 6\tabcolsep) * \real{0.2667}}
  >{\raggedleft\arraybackslash}p{(\linewidth - 6\tabcolsep) * \real{0.2667}}@{}}
\toprule\noalign{}
\begin{minipage}[b]{\linewidth}\raggedright
Reads after span
\end{minipage} & \begin{minipage}[b]{\linewidth}\raggedleft
Span insufficient
\end{minipage} & \begin{minipage}[b]{\linewidth}\raggedleft
Span sufficient
\end{minipage} & \begin{minipage}[b]{\linewidth}\raggedleft
Helper also delivered
\end{minipage} \\
\midrule\noalign{}
\endhead
\bottomrule\noalign{}
\endlastfoot
Untrained & 12/12 & 12/12 & 11/12 \\
Trained & 12/12 & 2/12 & 0/12 \\
\end{longtable}
}

The trained model went looking every time the span lacked the defect and
largely stopped when it did not; the untrained model read in all three
situations. Filtering to cases where reading was required, the trained
model read only once against the untrained model's 3.42 times. Both arms
successfully completed all tasks.

\subsection{WHEN: does the moment the diagnostic arrives change the
outcome?}\label{when-does-the-moment-the-diagnostic-arrives-change-the-outcome}

In the \emph{checker grid} the model was given a draft change that
passes its visible test but fails a held-out one, and asked to review it
and submit. We compared injecting a diagnostic after every edit, one at
revision, one that rejects a defective draft and asks for repair at
submission time, and a no-checker control.

Left to itself, the model rarely touched the draft so the
after-every-edit channel fired on only 1 of 12 defects and that arm
matched the control at 1 of 12 accepted. One-shot delivery at revision
reached 10 of 12, and the gate, which rejects a defective submission and
asks for repair, reached 11 of 12. Its single miss was a repair that
came back type-clean and still failed the held-out test.

{\def\LTcaptype{none} % do not increment counter
\begin{longtable}[]{@{}
  >{\raggedright\arraybackslash}p{(\linewidth - 8\tabcolsep) * \real{0.1579}}
  >{\raggedleft\arraybackslash}p{(\linewidth - 8\tabcolsep) * \real{0.2105}}
  >{\raggedleft\arraybackslash}p{(\linewidth - 8\tabcolsep) * \real{0.2105}}
  >{\raggedleft\arraybackslash}p{(\linewidth - 8\tabcolsep) * \real{0.2105}}
  >{\raggedleft\arraybackslash}p{(\linewidth - 8\tabcolsep) * \real{0.2105}}@{}}
\toprule\noalign{}
\begin{minipage}[b]{\linewidth}\raggedright
12 seeded defect/clean pairs
\end{minipage} & \begin{minipage}[b]{\linewidth}\raggedleft
Control
\end{minipage} & \begin{minipage}[b]{\linewidth}\raggedleft
After every edit
\end{minipage} & \begin{minipage}[b]{\linewidth}\raggedleft
One-shot
\end{minipage} & \begin{minipage}[b]{\linewidth}\raggedleft
Gate
\end{minipage} \\
\midrule\noalign{}
\endhead
\bottomrule\noalign{}
\endlastfoot
Bad completion accepted & 11/12 & 11/12 & 2/12 & 1/12 \\
Revision tokens, defect & 787 & 754 & 1,210 & 1,380 \\
Wall time, defect (s) & 27.8 & 14.5 & 68.3 & 62.7 \\
Revision tokens, clean & 591 & --- & 619 & 591 \\
Wall time, clean (s) & 9.4 & 9.9 & 9.6 & 9.5 \\
\end{longtable}
}

\emph{Mean edits 0.08 against 0.17 in control.}

Defects had to be seeded because capable models rarely left
checker-detectable ones given the relatively simple tasks we were
testing.

To test how the timing of diagnostics matters, the \emph{delivery
timing} suite put Qwen2.5-Coder-7B-Instruct and Pyrefly through six
delivery points. The arms were no feedback, batched at yield, batched
after each edit, and three variations of live channels that interrupt
mid-generation, splicing the diagnostic into the stream as Hooper et
al.~\citep{hooper2026speculative} do for asynchronous tool results. The
model was the stock instruct checkpoint in every arm.

{\def\LTcaptype{none} % do not increment counter
\begin{longtable}[]{@{}
  >{\raggedright\arraybackslash}p{(\linewidth - 6\tabcolsep) * \real{0.2308}}
  >{\raggedright\arraybackslash}p{(\linewidth - 6\tabcolsep) * \real{0.2308}}
  >{\raggedleft\arraybackslash}p{(\linewidth - 6\tabcolsep) * \real{0.3077}}
  >{\raggedright\arraybackslash}p{(\linewidth - 6\tabcolsep) * \real{0.2308}}@{}}
\toprule\noalign{}
\begin{minipage}[b]{\linewidth}\raggedright
Arm
\end{minipage} & \begin{minipage}[b]{\linewidth}\raggedright
When the diagnostic arrives
\end{minipage} & \begin{minipage}[b]{\linewidth}\raggedleft
Resolved
\end{minipage} & \begin{minipage}[b]{\linewidth}\raggedright
Difference vs no feedback
\end{minipage} \\
\midrule\noalign{}
\endhead
\bottomrule\noalign{}
\endlastfoot
none & never & 81/168 = 0.482 & --- \\
lazy & batched, at the model's next yield & 89/168 = 0.530 & +0.048
{[}−0.018, +0.113{]} \\
eager & batched, immediately after each edit & 88/168 = 0.524 & +0.042
{[}−0.024, +0.119{]} \\
naive & live mid-stream, told to fix each on arrival & 57/168 = 0.339 &
\textbf{−0.143 {[}−0.208, −0.077{]}} \\
plain & live mid-stream & 77/168 = 0.458 & −0.024 {[}−0.113,
+0.048{]} \\
gated & live mid-stream, only when the file parses & 81/168 = 0.482 &
+0.000 {[}−0.054, +0.054{]} \\
\end{longtable}
}

\emph{14 tasks x 12 seeds = 168 paired rollouts per arm. Differences are
task-clustered bootstraps; significance is paired exact McNemar under
Benjamini-Hochberg FDR 5\% across all 15 pairwise contrasts (cutoff
p\textless=0.0105). Only the naive arm's five contrasts clear it; among
the other five arms the smallest p is 0.119. Arms are \texttt{A},
\texttt{C-lazy}, \texttt{C-eager}, \texttt{D-naive}, \texttt{D-plain}
and \texttt{D-gate} in the run artifacts.}

Batching at a yield, batching after every edit and interrupting
mid-generation all land within noise of each other and of no feedback at
all. What hurt was the instruction. The naive live arm differs from the
plain live arm only by a sentence telling the model to fix each
diagnostic before moving on, and removing that sentence recovers most of
the deficit (+0.119 {[}+0.036, +0.202{]}, p=0.0105). Most ungated
diagnostics were about the model's own half-finished edit, so the
channel is noisy as described. It is the standing order to act on every
message, not the noise, that impacts the fix rate though.

\section{Conclusions, if you use a coding
agent}\label{conclusions-if-you-use-a-coding-agent}

LSP services can aid coding agents, but require more consideration than
simply making additional tools available.

\textbf{Keep type annotations correct and present.} Regardless of tools,
having correct, visible types improved how agents navigated the source.

\textbf{Run your type checker at the end of the agent's turn and make it
blocking.} This offered the clearest benefit and was the cheapest to
wire; it costs no tokens on clean work. Count how often it blocks, how
often the repair passes, and how often it rejects valid work.

\textbf{Make navigation tools available and encourage their use for
larger and more complex codebases.} For code in context navigation tools
add nothing, and for most well-typed code models will read and retrieve
correct spans from the appropriate source. Code navigation tools offer
the most gains where the target is somewhat obfuscated, as with
generated stubs, compiled extensions, vendored packages, or objects
assembled by a factory, registry or container.

\textbf{Do not assume a tool saves tokens unless it substitutes for
existing calls.} When adding a definition tool you can end up still
spending tokens on the ranged read: identify and measure duplicative
invocations. Prompting may help, but deep changes to this behavior could
require fine-tuning.

\section{Future work}\label{future-work}

\textbf{Retrieval calibration.} Agents retrieve constantly and mostly
indiscriminately: grep hits, file chunks, documentation, search results,
other tools' output. If ``do I have enough?'' is trainable, it is a
general lever on agent cost. We have seen tool-use post-training work
well, and enable the model to discriminate on tool calls. Can we train
to add additional code understanding channels, and to choose between
them intelligently?

\textbf{Tools co-designed with the agent.} The tools tested were built
for a human, who can ignore a transient diagnostic. An agent inherits
those outputs, but results suggest they lack some of the discrimination
in how to use them. Can a service better anticipate what an agent needs,
suppress likely-unhelpful messages and give more directive feedback on
actionable ones?

\textbf{Assessing a codebase in advance.} Results suggest types are
beneficial to coding agents regardless of the availability of LSP tools,
and adding LSP tooling offers wins primarily as codebases become more
complex. It is difficult to evaluate whether a given code base will
benefit from more aggressive usage of these types of tools however. Are
there measures we can use to assess this and predict the benefit, such
as annotation density, indirection depth between call and definition
sites, or the share of types resolving only through a stub or compiled
boundary?

\textbf{Which tasks benefit.} Everything here is dispatch ambiguity and
small seeded defects in readable workspaces. Semantic tooling may matter
in proportion to the sheer volume of context required, as seen in wide
refactors, cross-module renames, API migrations. Can we create tasks and
environments which evaluate these kinds of changes, and use them to
identify useful policies for LSP-related tools?

\section{Appendix: task construction and apparatus
checks}\label{appendix-task-construction-and-apparatus-checks}

Checker defects were screened so that each is coherent (the target file
parses and leaves no unimplemented stub), passes the visible test, fails
a held-out test, and carries a single semantic diagnostic in the target
file; each clean control is that defect's own validated gold. The
screens on the repository scan were environment, leakage, ambiguity and
discriminating lookup.

In the sufficiency experiment, a scripted adversary searched 17
one-parameter repair forms and found none that reached the held-out test
without retrieval, so the instances cannot be solved by guessing. The
harness also re-shows any file the agent has edited, which would have
delivered the hidden line without a read; we redacted that post-edit
view while leaving deliberate reads and grep available. The training set
behind that result contained no instances where a read was still
required after the span, so the selectivity the trained model shows was
never demonstrated to it; the result establishes suppression of the
post-span read rather than a learned conditional policy, and it was
tested against a single form of insufficiency.

The static AST resolver behind the definition operation was checked
against a live Pyrefly daemon at two levels: 12 of 12 agreement on the
synthetic task symbols, and the same outcome with byte-identical input
tokens on 22 paired real-library cells. That agreement covers definition
lookup on workspaces of the kind used here, not the rest of the
language-server surface. A further check on real library symbols agreed
9 of 11 with no disagreements, the two gaps being re-exports where
Pyrefly returned null and the resolver did not, so the two are not
interchangeable in general.

The shell-agent case study ran a single seed over 3 tasks with one model
and a 60-step cap, under which only 2 of 9 arms converged. Its counts
are a case series read alongside the controlled grids, not a matched
comparison with them.

%% paper/references.tex
%% Included after the document body (pandoc --include-after-body).
%% Emits the References section from the vendored references.bib.
%%
%% The References list now reflects REAL citations: report-filters.lua converts
%% the 8 related-work arXiv links into \citep{...} (keeping the visible work
%% name) and attaches \citep{ross2011dagger} at the "DAgger-style relabel"
%% anchor. Only those cited works are listed (9 entries). The remaining .bib
%% entries stay available for when more of the text is wired to \citep.
\bibliographystyle{plainnat}
\bibliography{references}

\end{document}
